\SourceLeaf{228}%
son has taken, and the results he has reached agree almost throughout with my own reflections, some of which I began 30—35 years ago. I am indeed astonished beyond measure at this. My intention was to let none of my own work become known in my lifetime. I have so far put little of it on paper. Most people have no proper sense of what matters here, and I have found only a few who took a particular interest in what I told them. To be able to do so, one must first have felt very keenly what is actually lacking, and most people have no clear idea of it. My intention, however, was in time to put everything on paper so that at least it would not perish with me one day.

I am thus greatly surprised that I can now be spared this effort, and it gives me the greatest pleasure that it is the son of my old friend who has anticipated me in so remarkable a way.

I find the notation very expressive and concise. I do think it would be good to give some of the principal concepts definite names, in addition to signs or letters, and I thought of several such names long ago. As long as one thinks the matter through in direct intuition alone, one needs no names or signs. They are needed only when one wants to communicate with others. Thus, e.~g., the surface your son calls $F$ could be called a parasphere, and the line $L$ a paracycle. In essence they are a spherical surface and a circle, each of infinite radius. One might call the aggregate of all points at an equal distance from a straight line in the same plane a hypercycle, and similarly use hypersphere. These are all trifling matters. It is the substance that matters, not the form.

In some parts of the investigation I took somewhat different paths. As a specimen I enclose a purely geometric proof, in outline, of the theorem that the difference between the sum of the angles of a triangle and $180^{0}$ is proportional to the area of the triangle.

\Needspace{6\baselineskip}
I. The aggregate of three straight lines $ab$, $cd$, $ef$ so situated that $ab\TripleParallel dc$, $cd\TripleParallel fe$, $ef\TripleParallel ba$ forms a figure that I call $T$. It can be
